\Needspace{12\baselineskip}
\originalchapter{热对流与Lorenz模型}
本章先使用Rayleigh--Bénard三模模型说明机制；原12.006J第20--21讲的环管热虹吸模型给b为1，完整环管推导与本章模型的逐系数缩放另见后文“黏性输运与低维模型的关键推导”。两种装置和b值分别保留。
\originalsection{传导态失稳}
自由流体的Rayleigh--Bénard模型由不可压缩动量、浮力和温度方程组成。以层厚、热扩散时间等作尺度，扰动温度 $\theta$ 相对于静止线性温度分布定义，垂直速度 $w$ 与温度满足线性方程
\[
Pr^{-1}\partial_t\nabla^2w=\nabla_h^2\theta+\nabla^4w,
\qquad \partial_t\theta=\nabla^2\theta+Ra\,w.
\]
对动量方程取两次旋度消除压力，再取竖向分量便得到第一式；不可压缩使压力梯度不需要单独求解。上下温度固定且应力自由时，$w=\theta=w_{zz}=0$ 在 $z=0,1$，正弦垂向模可以同时满足边界。

水平波数 $k$、垂向基模 $\pi$ 对应总 $q^2=k^2+\pi^2$。若 $w=A\cos(kx)\sin\pi z$、$\theta=B\cos(kx)\sin\pi z$，线性幅度矩阵为 $\dot A=-Pr\,q^2A+Pr(k^2/q^2)B$、$\dot B=Ra A-q^2B$。不稳定阈值 $Ra=q^6/k^2$，对 $k$ 最小化得 $k_c=\pi/\sqrt2$、$Ra_c=27\pi^4/4$。刚性边界不能用这个正弦函数同时满足全部条件，因此临界值不同。

\originalsection{三模截断的物理含义}
为描述对流超过阈值后的非线性，保留一个卷流幅度 $X$、一个与之同水平波数的温度幅度 $Y$、一个水平平均温度修正 $Z$。卷流与温度相乘产生垂向二次谐波，因此 $Z$ 不是随意加入的第三坐标。将场展开在相应三种三角函数上，投影到这些模，重新缩放时间和幅度后得到
\[
\dot X=Pr(Y-X),\quad \dot Y=rX-Y-XZ,\quad \dot Z=XY-bZ.
\]
$r=Ra/Ra_c$，对水平波数与垂向基模比值 $a=k/\pi$，三模模型的 $b=4/(1+a^2)$；取临界水平波数给 $b=8/3$。$X$ 表示循环强弱与方向，$Y$ 是左右温差，$Z$ 改变平均垂直温度剖面。$-XZ$ 限制温度驱动，$XY$ 产生平均温度修正。

截断丢弃了其他模，远离临界时没有一般误差保证。Lorenz模型的混沌是对非线性机制的示范，不等于此三变量轨迹可完整预测任何真实对流实验。

\originalsection{固定点和失稳}
传导固定点为零，对流固定点 $X=Y=\pm\sqrt{b(r-1)}$、$Z=r-1$，仅在 $r\ge1$ 存在。零点在 $r>1$ 出现正增长方向；非零点的三阶特征多项式为
\[
\lambda^3+(Pr+b+1)\lambda^2+b(Pr+r)\lambda+2Prb(r-1)=0.
\]
全部系数正时还需Routh--Hurwitz条件 $(Pr+b+1)b(Pr+r)>2Prb(r-1)$。若 $Pr>b+1$，失稳临界 $r_H=Pr(Pr+b+3)/(Pr-b-1)$；若分母非正，不存在同样的正临界表达式。这个参数条件在推导中必须保留。

负散度 $-(Pr+1+b)$ 保证相体积持续收缩，但在混沌吸引子上仍有一个伸长方向。卷流偶尔改变旋转方向，轨迹在两个对流分支附近交替停留；外观不规则不自动等于已经测得正Lyapunov指数。

\originalsection{Rayleigh数与模态的符号约定}
温度场 $T=T_{
m cond}(z)+\theta$，基本梯度 $T'_{
m cond}<0$，所以扰动热方程含 $-wT'_{
m cond}>0$，上升扰动搬运更暖流体。压力从动量的两次旋度消去产生垂向波数因子，热与速度模必须满足同一边界条件。正文将 $Ra$ 放在热方程中是温度幅度的缩放选择，等价于常见把它放在浮力方程；不可从两种约定各取半式拼接。

对垂向第 $n$ 模，临界曲线 $Ra_n(k)=(k^2+n^2\pi^2)^3/k^2$；最先为 $n=1$，最小化过程给 $k_c^2=\pi^2/2$。波长 $2\pi/k_c$ 相对层厚是可观测的空间尺度，临界 $Ra$ 是失稳阈值，二者不是同一个输出。

\originalsection{投影为何产生Lorenz的三个非线性项}
二维卷流可用流函数，选 $\psi\propto X\sin(kx)\sin(\pi z)$，温度选 $Y\cos(kx)\sin(\pi z)-Z\sin(2\pi z)$。对流项由Jacobian $J(\psi,\theta)=\psi_x\theta_z-\psi_z\theta_x$ 给出：第一温度模与卷流相乘生成水平平均 $\sin(2\pi z)$，对应 $XY$；平均温度修正与卷流又影响第一模，对应 $-XZ$。投影的正交积分删去不在三模中的部分，并重新缩放系数为标准形式。

这种闭合选择有物理动机却不是严格不变子空间；被删除模仍可被非线性激发。越近传导失稳，少数慢模通常越有价值，远超阈值则应扩展模型或回到场方程。课程以它展示混沌的发现，不意味着三变量足以完成环境流体预测。

\originalsection{固定点的稳定性与Hopf条件}
零点的 $Z$ 模衰减率 $-b$，$X,Y$ 二维块行列式 $Pr(1-r)$，所以 $r>1$ 有一正方向。非零点处三阶多项式的系数可用Jacobian直接展开，Routh--Hurwitz要求 $a_1,a_2,a_3>0,a_1a_2>a_3$。设置后一条件等号得到正文 $r_H$；正临界存在要求 $Pr>b+1$。

对 $Pr=10,b=8/3$，$r_H\simeq24.73684$。Hopf临界本身并不保证紧接着产生稳定小环；Lorenz系统还可能有亚临界与全局轨道结构。$r=28$ 的常用混沌参数不能仅从“28大于24.7”证明混沌，必须进一步观测长期轨迹/切向指数和回返结构。

\originalsection{吸收域、奇异吸引子与计算证据}
负散度只给体积收缩，未单独保证轨迹有界。取加权二次函数 $V=rX^2+PrY^2+Pr(Z-2r)^2$，求导中 $XYZ$ 项与交叉项相消，得到
\[
\frac12\dot V=-Pr[rX^2+Y^2+b(Z-r)^2-br^2].
\]
椭球外足够大时负，因而可构造吸收域；这一步把“有限范围吸引”与“体积耗散”区分。$r=0$ 等特殊参数可用其他正定权重，本文关注 $r>0$。

附带程序实际以DOP853积分 $r=10,28$、初值 $(.01,.01,.01)$、总时长200，保存积分容差与采样输出。$r=10$ 收敛到对流平衡 $X=Y\simeq4.89898,Z=9$；$r=28$ 后半段有多次卷流换向，局部非零固定点特征值一实负和一对实部约 $.09396$ 的复根。图显示模型轨迹，未假装来自实体对流测量。对混沌，逐点长时轨迹随容差的差别并非普通函数收敛，应比较统计和切向增长。

\originalsection{R\"ossler：伸展和折返的另一种机制}
12.006J第19讲实际以R\"ossler为例，原插图有第三方许可例外，本节据方程独立推导并用自算轨迹绘图。模型
\[
\dot x=-y-z,\qquad\dot y=x+ay,\qquad\dot z=b+z(x-c)
\]
在 $z$ 小时近为 $\ddot x-a\dot x+x=0$；$a>0$ 是负阻尼，使平面轨迹向外螺旋。当 $x$ 增到 $c$ 附近，$z$ 的增长项增强，随后 $-z$ 把 $x$ 拉回，形成折返。$b>0$ 时 $z$ 近零也有输入，不能严格把 $x<c$ 都说成 $\dot z<0$，那只是 $b$ 小而 $z(c-x)>b$ 的区间判断。

固定点满足 $x=az,y=-z,az^2-cz+b=0$，所以 $c^2\ge4ab$ 时有两根。散度 $a+x-c$ 随位置变化，不能像Lorenz那样由一个恒负常数判断所有区域。实际轨迹上的时间平均散度可核验总体收缩。

本卷明示选择 $a=b=.2,c=5.7$、初值 $(1,0,0)$，积分并沿切向方程反复归一化。图是模型轨迹投影；QA给有限时间最大Lyapunov估计、时长和容差，不把有限时间估计当作极限存在的严格证明。它说明不规则吸引可以来自伸展/折返/总体收缩，不依赖Lorenz那样在两固定点附近交换旋转。
\begin{center}\begin{tikzpicture}\begin{axis}[width=.7\linewidth,height=5cm,xlabel={$x$},ylabel={$y$}]
\addplot[main,thin] table[x=x,y=y,col sep=comma]{data/rossler-plot.csv};
\end{axis}\end{tikzpicture}\end{center}
